\documentclass{article}
\usepackage{amsfonts}
\usepackage{amsmath}
\usepackage{amsthm}
\usepackage{amssymb}
\usepackage{mathtools}% http://ctan.org/pkg/mathtools
\usepackage{hyperref}
\usepackage{caption}
\usepackage{nameref}
\DeclarePairedDelimiter\ceil{\lceil}{\rceil}
\DeclarePairedDelimiter\floor{\lfloor}{\rfloor}
\linespread{1.3} % Roughly 1.5 spacing
\usepackage[margin=0.75in]{geometry}
\usepackage[fontsize=11pt]{scrextend}

\title{MAT 324 Preparation}
\date{}

\begin{document}



\section*{Exercises 1A}
\textbf{1. } Suppose $f:[a,b]\rightarrow \mathbb{R}$ is a bounded function such that $L(f,P,[a,b])=U(f,P,[a,b])$ for some partition $P$ of $[a,b]$. Prove that $f$ is the constant function on $[a,b]$.
\newline
\newline
Let $a=x_0,\ldots,x_n=b$ denote the partition $P$ of $[a,b]$. Since the lower and upper Riemann sums are equal, we have
\begin{align*}
    \sum_{j=1}^n(x_j-x_{j-1})\underset{[x_{j-1},x_j]}{\sup}f=\sum_{j=1}^n(x_j-x_{j-1})\underset{[x_{j-1},x_j]}{\inf}f.\\
    \implies
    \sum_{j=1}^n(x_j-x_{j-1})\Big(\underset{[x_{j-1},x_j]}{\sup}f-\underset{[x_{j-1},x_j]}{\inf}f\Big)=0\\
\end{align*}
The length of each interval is strictly positive, and $\underset{[x_{j-1},x_j]}{\sup}f\geq \underset{[x_{j-1},x_j]}{\inf}f$ for each $j$, so the above sum consists of only non-negative terms. Thus, we have $(x_j-x_{j-1})(\underset{[x_{j-1},x_j]}{\sup}f-\underset{[x_{j-1},x_j]}{\inf}f)=0$ for each interval, and $(\underset{[x_{j-1},x_j]}{\sup}f=\underset{[x_{j-1},x_j]}{\inf}f)$. Thus, $f$ is constant on every interval $[x_{j-1},x_j]$.
\\
\newline
\newline
$f(x)=f(a)$ for every $x\in[x_1,x_2]$. Suppose $f$ is $f(a)$ on an interval $[x_{j-1},x_j]$. Since $f$ is also constant on the interval $[x_j,x_{j+1}]$, $f(x)=f(x_j)=f(a)$ for all $x\in [x_j,x_{j+1}]$. By mathematical induction, $f(x)=f(a)$ on every interval $[x_j,x_{j+1}]$. Therefore, $f$ is constant on the entire domain.\qed
\newline
\newline
\textbf{2. } Suppose that $a\leq s < t \leq b$. Define $f:[a,b]\rightarrow \mathbb{R}$ by 
\begin{equation*}
    f(x)=
    \begin{cases*}
        1 &\text{if }$s<x<t$,\\
        0 &\text{otherwise.}
    \end{cases*}
\end{equation*}
Prove that $f$ is Riemann integrable on $[a,b]$ and that $\int_{a}^{b}f=t-s$. 
\newline
\newline
Consider the partition $P$ of $[a,b]$ into $n\in\mathbb{N}$ intervals of equal length given by $x_i=a+i\frac{b-a}n$. Notice that 
\begin{equation*}
    L(f,P,[a,b])=\sum_{i=1}^{n}d\underset{[x_{i-1},x_i]}{\inf}f\geq(\floor{\frac{n(t-s)}{b-a}}-2)\frac{b-a}{n},
\end{equation*}
and similarly
\begin{equation*}
    U(f,P,[a,b])=\sum_{i=1}^{n}d\underset{[x_{i-1},x_i]}{\sup}f\leq\floor{\frac{n(t-s)}{b-a}}\frac{b-a}{n}.
\end{equation*}
Then,
\begin{equation*}
    L(f,[a,b])
    \geq \underset{n\in \mathbb{N}}{\sup}(\floor{\frac{n(t-s)}{b-a}}-2)\frac{b-a}{n}
    = t-s
    = \underset{n\in \mathbb{N}}{\inf}(\floor{\frac{n(t-s)}{b-a}})\frac{b-a}{n}
    \geq U(f,[a,b]).
\end{equation*}
But, since $L(f,[a,b])\leq U(f,[a,b])$ we have $L(f,[a,b])= U(f,[a,b])=t-s$, and hence $\int_{a}^{b}f=t-s$.\qed
\newpage\noindent
\textbf{3. }Suppose $f:[a,b]\rightarrow \mathbb{R}$ is a bounded function. Prove that $f$ is Riemann integrable if and only if for each $\epsilon >0$ there exists a partition $P$ of $[a,b]$ such that 
\begin{equation*}
    U(f,P,[a,b])-L(f,P,[a,b])<\epsilon.
\end{equation*}
($\implies$) Suppose $f$ is Riemann integrable, and let $I:=\int_{a}^{b}f$. Let $\epsilon >0$. Since $I=\underset{P}\sup L(f,P,[a,b])=\underset{P}\inf U(f,P,[a,b])$, there is a partition $P_1$ so that $L(f,P_1,[a,b])>I-\frac\epsilon 2$ and a partition $P_2$ so that $U(f,P_2,[a,b])<I+\frac{\epsilon}2$. Let $P$ be the partition obtained by taking the union of $P_1$ and $P_2$. Since $P$ is a finer partition than both $P_1$ and $P_2$, we have
\begin{equation*}
    L(f,P,[a,b])\geq L(f,P_1,[a,b])>I-\epsilon/2
\end{equation*}
and
\begin{equation*}
    U(f,P,[a,b])\leq U(f,P_2,[a,b])<I+\epsilon/2.
\end{equation*}
Hence, $|L(f,P[a,b])-I|<\epsilon/2$ and $|U(f,P[a,b])-I|<\epsilon/2$. By the triangle inequality, we have 
\begin{align*}
    U(f,P,[a,b])-L(f,P,[a,b])&=\leq|U(f,P,[a,b])-L(f,P,[a,b])|\\
    &=|U(f,P,[a,b])-I+I-L(f,P,[a,b])|\\
    &\leq|U(f,P,[a,b])-I|+|I-L(f,P,[a,b])|\\
    &<2(\epsilon/2)=\epsilon,
\end{align*}
as desired.
\newline
\newline
($\impliedby$) Suppose that $f$ is not Riemann integrable. Let $L:=L(f,[a,b])$ and $U:=U(f,[a,b])$ be the lower and upper Riemann integrals, where $L\neq U$. Since $L(f,P,[a,b])\leq U(f,P,[a,b])$ for any partition, $L\leq U$. Let $\epsilon:=\frac{U-L}2$. Then, for any partition $P$ of $[a,b]$, we have 
\begin{align*}
    U(f,P,[a,b])-L(f,P,[a,b])&=\geq \underset{P}{\inf}U(f,P,[a,b])-\underset{P}{\sup}L(f,P,[a,b])\\
    &= U-L\\
    &>\epsilon.
\end{align*}
\qed
\newpage \noindent
\textbf{4. }Suppose $f,g:[a,b]\rightarrow \mathbb{R}$ are Riemann integrable. Prove that $f+g$ is Riemann integrable and that $\int_{a}^{b}(f+g)=\int_{a}^{b}f+\int_{a}^{b}g$. 
\newline
\newline
Recall that the infimum and supremum distribute over the minkowski sums of sets. Then, we have
\begin{align*}
    L(f+g,[a,b])&=\underset{P}{\sup} L(f+g,P,[a,b])\\
    &=\underset{P}{\sup} \sum_{j=1}^{n} (x_j-x_{j-1})\underset{[x_{j-1},x_j]}{\inf} f+g\\
    &=\underset{P}{\sup} \sum_{j=1}^{n} (x_j-x_{j-1})\underset{[x_{j-1},x_j]}{\inf} f +\underset{P}{\sup} \sum_{j=1}^{n} (x_j-x_{j-1})\underset{[x_{j-1},x_j]}{\inf} f\\
    &= \int_{a}^{b}f+\int_{a}^{b}g\\
    &= \underset{P}{\inf} \sum_{j=1}^{n} (x_j-x_{j-1})\underset{[x_{j-1},x_j]}{\sup} f +\underset{P}{\inf} \sum_{j=1}^{n} (x_j-x_{j-1})\underset{[x_{j-1},x_j]}{\sup} f\\
    &=\underset{P}{\inf} \sum_{j=1}^{n} (x_j-x_{j-1})\underset{[x_{j-1},x_j]}{\sup} f+g\\
    &=\underset{P}{\inf} U(f+g,P,[a,b])\\
    &=U(f+g,[a,b])\\
\end{align*}
Hence, $L(f+g,[a,b])=U(f+g,[a,b])=\int_{a}^{b}f+\int_{a}^{b}g$. \qed

\end{document}
